\section{Parameter Estimation and Optimization}
\subsection{Individual-level maximum likelihood}

Each participant was fitted independently. For model parameters \(\vartheta\), estimation solves
\begin{equation}
    \hat{\vartheta}=\arg\min_{\vartheta}\NLL(\vartheta).
\end{equation}
The fitted parameters are therefore those that make that participant's observed sequence of choices most probable under the model.

\subsection{Bounded optimization with L-BFGS-B}

The fitting problems contain continuous parameters, a smooth likelihood objective, explicit parameter bounds, and relatively few dimensions. L-BFGS-B was used because it is an efficient gradient-based optimizer designed for continuous optimization with box constraints~\cite{byrd1995}. Because it is a local optimizer, however, its solution can depend on the starting point.

\subsection{Multiple starting points}

To reduce sensitivity to local minima, each participant/model fit used multiple starts. Every start was optimized independently and the valid solution with the lowest NLL was retained. This is methodologically important because a poor local optimum can change the apparent model comparison.

For \QL{}, only two parameters, \((\eta,\beta)\), are fitted. The objective was evaluated on a structured two-dimensional grid, distinct promising local-minimum regions were identified, and up to 10 grid-local minima were used as L-BFGS-B starting points. Because the parameter space is only two-dimensional, this provides a relatively systematic search rather than relying on arbitrary random starts.

For \PVL{}, the four-dimensional parameter space \((\phi,\alpha,\lambda,c)\) makes a dense Cartesian grid less practical. The analysis therefore used 32 scrambled Sobol points with fixed seed 42. Sobol sequences provide space-filling, low-discrepancy coverage of a bounded multidimensional region~\cite{sobol1967}.

\subsection{Parameter bounds and fixed settings}

\begin{table}[H]
    \centering
    \caption{Primary fitting settings used in the final analysis.}\label{tab:fitting-settings}
    \small
    \begin{tabular}{@{}
        >{\RaggedRight\arraybackslash}p{0.30\textwidth}
        >{\RaggedRight\arraybackslash}p{0.24\textwidth}
        >{\RaggedRight\arraybackslash}p{0.38\textwidth}
        @{}}
        \toprule
        Quantity                        & Final setting              & Rationale                                              \\
        \midrule
        Q learning rate \(\eta\)        & \([0,1]\)                  & Standard delta-rule range                              \\
        Q inverse temperature \(\beta\) & \([0,100]\)                & Upper cap selected after sensitivity analysis          \\
        Payoff scale                    & 100                        & Outcome rescaling for the fitted value scale           \\
        PVL learning rate \(\phi\)      & \([10^{-6},1-10^{-6}]\)    & Numerical approximation to the open interval \((0,1)\) \\
        Outcome sensitivity \(\alpha\)  & \([10^{-6},2]\)            & Expanded IGT/PVL range used in later modeling~\cite{haines2018} \\
        Loss aversion \(\lambda\)       & \([10^{-6},10]\)           & Expanded IGT/PVL range used in later modeling~\cite{haines2018} \\
        Response consistency \(c\)      & \([0,5]\)                  & PVL response-consistency range~\cite{ahn2008}          \\
        Q initialization                & Up to 10 \mbox{grid-local minima} & Structured search in two dimensions                    \\
        PVL initialization              & 32 scrambled Sobol points  & Space-filling search in four dimensions                \\
        Sobol seed                      & 42                         & Reproducibility                                        \\
        \bottomrule
    \end{tabular}
\end{table}
